\section{Results}\label{sec:results}
% Every number and figure here should come from scripts/ so it can be regenerated.
% Generated by scripts/meshing_results.py from data/meshing_ext8.csv; do not edit.
\newcommand{\MeshGeometriesHuman}{50}
\newcommand{\MeshGeometriesPctHuman}{100}
\newcommand{\MeshGeometriesPorcine}{44}
\newcommand{\MeshGeometriesPctPorcine}{100}
\newcommand{\MeshGeometriesAll}{94}
\newcommand{\MeshGeometriesPctAll}{100}
\newcommand{\MeshMeshedHuman}{41}
\newcommand{\MeshMeshedPctHuman}{82}
\newcommand{\MeshMeshedPorcine}{34}
\newcommand{\MeshMeshedPctPorcine}{77}
\newcommand{\MeshMeshedAll}{75}
\newcommand{\MeshMeshedPctAll}{80}
\newcommand{\MeshCleanHuman}{32}
\newcommand{\MeshCleanPctHuman}{64}
\newcommand{\MeshCleanPorcine}{22}
\newcommand{\MeshCleanPctPorcine}{50}
\newcommand{\MeshCleanAll}{54}
\newcommand{\MeshCleanPctAll}{57}
\newcommand{\MeshOrientationHuman}{6}
\newcommand{\MeshOrientationPctHuman}{12}
\newcommand{\MeshOrientationPorcine}{3}
\newcommand{\MeshOrientationPctPorcine}{7}
\newcommand{\MeshOrientationAll}{9}
\newcommand{\MeshOrientationPctAll}{10}
\newcommand{\MeshNegativeHuman}{3}
\newcommand{\MeshNegativePctHuman}{6}
\newcommand{\MeshNegativePorcine}{9}
\newcommand{\MeshNegativePctPorcine}{20}
\newcommand{\MeshNegativeAll}{12}
\newcommand{\MeshNegativePctAll}{13}
\newcommand{\MeshFailedHuman}{9}
\newcommand{\MeshFailedPctHuman}{18}
\newcommand{\MeshFailedPorcine}{10}
\newcommand{\MeshFailedPctPorcine}{23}
\newcommand{\MeshFailedAll}{19}
\newcommand{\MeshFailedPctAll}{20}
\newcommand{\MeshOutletClosedHuman}{7}
\newcommand{\MeshOutletClosedPctHuman}{14}
\newcommand{\MeshOutletClosedPorcine}{3}
\newcommand{\MeshOutletClosedPctPorcine}{7}
\newcommand{\MeshOutletClosedAll}{10}
\newcommand{\MeshOutletClosedPctAll}{11}
\newcommand{\MeshTouchingOpeningsHuman}{2}
\newcommand{\MeshTouchingOpeningsPctHuman}{4}
\newcommand{\MeshTouchingOpeningsPorcine}{1}
\newcommand{\MeshTouchingOpeningsPctPorcine}{2}
\newcommand{\MeshTouchingOpeningsAll}{3}
\newcommand{\MeshTouchingOpeningsPctAll}{3}
\newcommand{\MeshNonManifoldWallHuman}{0}
\newcommand{\MeshNonManifoldWallPctHuman}{0}
\newcommand{\MeshNonManifoldWallPorcine}{3}
\newcommand{\MeshNonManifoldWallPctPorcine}{7}
\newcommand{\MeshNonManifoldWallAll}{3}
\newcommand{\MeshNonManifoldWallPctAll}{3}
\newcommand{\MeshDuplicatedPatchHuman}{0}
\newcommand{\MeshDuplicatedPatchPctHuman}{0}
\newcommand{\MeshDuplicatedPatchPorcine}{3}
\newcommand{\MeshDuplicatedPatchPctPorcine}{7}
\newcommand{\MeshDuplicatedPatchAll}{3}
\newcommand{\MeshDuplicatedPatchPctAll}{3}
\newcommand{\MeshScanDefectsHuman}{2}
\newcommand{\MeshScanDefectsPorcine}{7}
\newcommand{\MeshScanDefectsAll}{9}
\newcommand{\FluidCellsMedianHuman}{364}
\newcommand{\FluidCellsQOneHuman}{312}
\newcommand{\FluidCellsQThreeHuman}{413}
\newcommand{\FluidCellsMinHuman}{247}
\newcommand{\FluidCellsMaxHuman}{543}
\newcommand{\FluidCellsMedianPorcine}{209}
\newcommand{\FluidCellsQOnePorcine}{190}
\newcommand{\FluidCellsQThreePorcine}{242}
\newcommand{\FluidCellsMinPorcine}{125}
\newcommand{\FluidCellsMaxPorcine}{437}
\newcommand{\FluidCellsMedianAll}{297}
\newcommand{\FluidCellsQOneAll}{225}
\newcommand{\FluidCellsQThreeAll}{390}
\newcommand{\FluidCellsMinAll}{125}
\newcommand{\FluidCellsMaxAll}{543}
\newcommand{\SolidCellsMedianHuman}{191}
\newcommand{\SolidCellsQOneHuman}{174}
\newcommand{\SolidCellsQThreeHuman}{239}
\newcommand{\SolidCellsMinHuman}{123}
\newcommand{\SolidCellsMaxHuman}{272}
\newcommand{\SolidCellsMedianPorcine}{124}
\newcommand{\SolidCellsQOnePorcine}{109}
\newcommand{\SolidCellsQThreePorcine}{138}
\newcommand{\SolidCellsMinPorcine}{94}
\newcommand{\SolidCellsMaxPorcine}{186}
\newcommand{\SolidCellsMedianAll}{165}
\newcommand{\SolidCellsQOneAll}{126}
\newcommand{\SolidCellsQThreeAll}{197}
\newcommand{\SolidCellsMinAll}{94}
\newcommand{\SolidCellsMaxAll}{272}
\newcommand{\FluidNonOrthoMedianHuman}{62}
\newcommand{\FluidNonOrthoQOneHuman}{60}
\newcommand{\FluidNonOrthoQThreeHuman}{63}
\newcommand{\FluidNonOrthoMinHuman}{52}
\newcommand{\FluidNonOrthoMaxHuman}{70}
\newcommand{\FluidNonOrthoMedianPorcine}{62}
\newcommand{\FluidNonOrthoQOnePorcine}{59}
\newcommand{\FluidNonOrthoQThreePorcine}{64}
\newcommand{\FluidNonOrthoMinPorcine}{55}
\newcommand{\FluidNonOrthoMaxPorcine}{70}
\newcommand{\FluidNonOrthoMedianAll}{62}
\newcommand{\FluidNonOrthoQOneAll}{59}
\newcommand{\FluidNonOrthoQThreeAll}{63}
\newcommand{\FluidNonOrthoMinAll}{52}
\newcommand{\FluidNonOrthoMaxAll}{70}
\newcommand{\FluidSkewMedianHuman}{1.31}
\newcommand{\FluidSkewQOneHuman}{1.13}
\newcommand{\FluidSkewQThreeHuman}{1.52}
\newcommand{\FluidSkewMinHuman}{0.67}
\newcommand{\FluidSkewMaxHuman}{3.27}
\newcommand{\FluidSkewMedianPorcine}{1.18}
\newcommand{\FluidSkewQOnePorcine}{0.90}
\newcommand{\FluidSkewQThreePorcine}{1.48}
\newcommand{\FluidSkewMinPorcine}{0.70}
\newcommand{\FluidSkewMaxPorcine}{2.64}
\newcommand{\FluidSkewMedianAll}{1.27}
\newcommand{\FluidSkewQOneAll}{1.05}
\newcommand{\FluidSkewQThreeAll}{1.49}
\newcommand{\FluidSkewMinAll}{0.67}
\newcommand{\FluidSkewMaxAll}{3.27}
\newcommand{\SolidNonOrthoMedianHuman}{67}
\newcommand{\SolidNonOrthoQOneHuman}{63}
\newcommand{\SolidNonOrthoQThreeHuman}{70}
\newcommand{\SolidNonOrthoMinHuman}{46}
\newcommand{\SolidNonOrthoMaxHuman}{75}
\newcommand{\SolidNonOrthoMedianPorcine}{68}
\newcommand{\SolidNonOrthoQOnePorcine}{66}
\newcommand{\SolidNonOrthoQThreePorcine}{76}
\newcommand{\SolidNonOrthoMinPorcine}{56}
\newcommand{\SolidNonOrthoMaxPorcine}{82}
\newcommand{\SolidNonOrthoMedianAll}{68}
\newcommand{\SolidNonOrthoQOneAll}{63}
\newcommand{\SolidNonOrthoQThreeAll}{71}
\newcommand{\SolidNonOrthoMinAll}{46}
\newcommand{\SolidNonOrthoMaxAll}{82}
\newcommand{\SolidSkewMedianHuman}{1.99}
\newcommand{\SolidSkewQOneHuman}{1.47}
\newcommand{\SolidSkewQThreeHuman}{2.42}
\newcommand{\SolidSkewMinHuman}{0.79}
\newcommand{\SolidSkewMaxHuman}{3.42}
\newcommand{\SolidSkewMedianPorcine}{1.38}
\newcommand{\SolidSkewQOnePorcine}{0.91}
\newcommand{\SolidSkewQThreePorcine}{1.98}
\newcommand{\SolidSkewMinPorcine}{0.70}
\newcommand{\SolidSkewMaxPorcine}{3.54}
\newcommand{\SolidSkewMedianAll}{1.74}
\newcommand{\SolidSkewQOneAll}{1.23}
\newcommand{\SolidSkewQThreeAll}{2.40}
\newcommand{\SolidSkewMinAll}{0.70}
\newcommand{\SolidSkewMaxAll}{3.54}
\newcommand{\TaubinMovedMedianHuman}{99.7}
\newcommand{\TaubinMovedQOneHuman}{99.5}
\newcommand{\TaubinMovedQThreeHuman}{99.7}
\newcommand{\TaubinMovedMinHuman}{99.0}
\newcommand{\TaubinMovedMaxHuman}{99.9}
\newcommand{\TaubinMovedMedianPorcine}{98.7}
\newcommand{\TaubinMovedQOnePorcine}{98.2}
\newcommand{\TaubinMovedQThreePorcine}{99.2}
\newcommand{\TaubinMovedMinPorcine}{97.4}
\newcommand{\TaubinMovedMaxPorcine}{99.9}
\newcommand{\TaubinMovedMedianAll}{99.6}
\newcommand{\TaubinMovedQOneAll}{98.9}
\newcommand{\TaubinMovedQThreeAll}{99.7}
\newcommand{\TaubinMovedMinAll}{97.4}
\newcommand{\TaubinMovedMaxAll}{99.9}
\newcommand{\SaddleMovedMedianHuman}{13.0}
\newcommand{\SaddleMovedQOneHuman}{11.3}
\newcommand{\SaddleMovedQThreeHuman}{15.0}
\newcommand{\SaddleMovedMinHuman}{9.0}
\newcommand{\SaddleMovedMaxHuman}{20.2}
\newcommand{\SaddleMovedMedianPorcine}{11.5}
\newcommand{\SaddleMovedQOnePorcine}{9.3}
\newcommand{\SaddleMovedQThreePorcine}{13.5}
\newcommand{\SaddleMovedMinPorcine}{4.8}
\newcommand{\SaddleMovedMaxPorcine}{18.1}
\newcommand{\SaddleMovedMedianAll}{12.0}
\newcommand{\SaddleMovedQOneAll}{10.0}
\newcommand{\SaddleMovedQThreeAll}{14.4}
\newcommand{\SaddleMovedMinAll}{4.8}
\newcommand{\SaddleMovedMaxAll}{20.2}
\newcommand{\NormalSmoothedMedianHuman}{0.7}
\newcommand{\NormalSmoothedQOneHuman}{0.2}
\newcommand{\NormalSmoothedQThreeHuman}{1.2}
\newcommand{\NormalSmoothedMinHuman}{0.0}
\newcommand{\NormalSmoothedMaxHuman}{2.4}
\newcommand{\NormalSmoothedMedianPorcine}{0.5}
\newcommand{\NormalSmoothedQOnePorcine}{0.3}
\newcommand{\NormalSmoothedQThreePorcine}{1.5}
\newcommand{\NormalSmoothedMinPorcine}{0.0}
\newcommand{\NormalSmoothedMaxPorcine}{9.2}
\newcommand{\NormalSmoothedMedianAll}{0.7}
\newcommand{\NormalSmoothedQOneAll}{0.2}
\newcommand{\NormalSmoothedQThreeAll}{1.3}
\newcommand{\NormalSmoothedMinAll}{0.0}
\newcommand{\NormalSmoothedMaxAll}{9.2}
\newcommand{\VolTaubinMedianHuman}{0.52}
\newcommand{\VolTaubinQOneHuman}{0.49}
\newcommand{\VolTaubinQThreeHuman}{0.55}
\newcommand{\VolTaubinMinHuman}{0.22}
\newcommand{\VolTaubinMaxHuman}{0.62}
\newcommand{\VolTaubinMedianPorcine}{0.58}
\newcommand{\VolTaubinQOnePorcine}{0.56}
\newcommand{\VolTaubinQThreePorcine}{0.60}
\newcommand{\VolTaubinMinPorcine}{0.42}
\newcommand{\VolTaubinMaxPorcine}{0.61}
\newcommand{\VolTaubinMedianAll}{0.55}
\newcommand{\VolTaubinQOneAll}{0.51}
\newcommand{\VolTaubinQThreeAll}{0.59}
\newcommand{\VolTaubinMinAll}{0.22}
\newcommand{\VolTaubinMaxAll}{0.62}
\newcommand{\VolSaddleMedianHuman}{0.11}
\newcommand{\VolSaddleQOneHuman}{0.05}
\newcommand{\VolSaddleQThreeHuman}{0.20}
\newcommand{\VolSaddleMinHuman}{$-$0.17}
\newcommand{\VolSaddleMaxHuman}{0.39}
\newcommand{\VolSaddleMedianPorcine}{0.11}
\newcommand{\VolSaddleQOnePorcine}{0.08}
\newcommand{\VolSaddleQThreePorcine}{0.18}
\newcommand{\VolSaddleMinPorcine}{$-$0.32}
\newcommand{\VolSaddleMaxPorcine}{0.30}
\newcommand{\VolSaddleMedianAll}{0.11}
\newcommand{\VolSaddleQOneAll}{0.06}
\newcommand{\VolSaddleQThreeAll}{0.19}
\newcommand{\VolSaddleMinAll}{$-$0.32}
\newcommand{\VolSaddleMaxAll}{0.39}
\newcommand{\VolTotalMedianHuman}{0.61}
\newcommand{\VolTotalQOneHuman}{0.55}
\newcommand{\VolTotalQThreeHuman}{0.75}
\newcommand{\VolTotalMinHuman}{0.13}
\newcommand{\VolTotalMaxHuman}{0.98}
\newcommand{\VolTotalMedianPorcine}{0.70}
\newcommand{\VolTotalQOnePorcine}{0.66}
\newcommand{\VolTotalQThreePorcine}{0.76}
\newcommand{\VolTotalMinPorcine}{0.12}
\newcommand{\VolTotalMaxPorcine}{0.87}
\newcommand{\VolTotalMedianAll}{0.68}
\newcommand{\VolTotalQOneAll}{0.58}
\newcommand{\VolTotalQThreeAll}{0.76}
\newcommand{\VolTotalMinAll}{0.12}
\newcommand{\VolTotalMaxAll}{0.98}
\newcommand{\TaubinDispMedianAll}{0.15}
\newcommand{\SaddleDispMedianAll}{0.35}
\newcommand{\SaddleDispMaxAll}{4.2}
\newcommand{\MeshNegativeCellsMax}{44}
\newcommand{\MeshNegativeCellsTotal}{173}
\newcommand{\MeshRuntimeMedianMin}{2}
\newcommand{\MeshRuntimeMaxMin}{8}
\newcommand{\MeshOpeningsMedianHuman}{16}
\newcommand{\MeshOpeningsMedianPorcine}{12}
% numbers quoted below, generated by scripts/meshing_results.py
\subsection{Meshing pipeline robustness}\label{sec:res-pipeline}
% Generated by scripts/meshing_results.py from data/meshing_ext8.csv; do not edit.
\begin{table}[tbp]
\centering
\footnotesize
\setlength{\tabcolsep}{3pt}
\caption{Outcome of the automated meshing workflow at $\varepsilon = 8$ for every geometry, and, for the geometries whose solid mesh passed the standard OpenFOAM mesh checks, the mesh size (median [interquartile range]), mesh quality (median [worst]) and the changes made to the surface and the extrusion (median [range]). A point counts as moved when it is displaced by more than \SI{0.02}{\milli\metre}; smoothing is the windowed-sinc smoothing of the surface, and volume changes are of the lumen.}
\label{tab:meshing}
\begin{tabular}{@{}lccc@{}}
\toprule
 & Human & Porcine & All \\
\midrule
\multicolumn{4}{l}{\textit{Outcome (geometries)}} \\
Geometries & 50 & 44 & 94 \\
Volume meshes generated & 41 & 34 & 75 \\
\quad Solid passes mesh checks & 32 & 22 & 54 \\
\quad Oriented-face errors only & 6 & 3 & 9 \\
\quad Negative-volume cells & 3 & 9 & 12 \\
No volume meshes & 9 & 10 & 19 \\
\quad Small outlet closed & 7 & 3 & 10 \\
\quad Openings touching (scan) & 2 & 1 & 3 \\
\quad Non-manifold wall (scan) & 0 & 3 & 3 \\
\quad Duplicated wall patch (scan) & 0 & 3 & 3 \\
\addlinespace
\multicolumn{4}{l}{\textit{Mesh size ($\times 10^{3}$ cells)}} \\
Fluid cells & 364 [312, 413] & 209 [190, 242] & 297 [225, 390] \\
Solid cells & 191 [174, 239] & 124 [109, 138] & 165 [126, 197] \\
\addlinespace
\multicolumn{4}{l}{\textit{Mesh quality}} \\
Fluid max.\ non-orthogonality (\si{\degree}) & 62 [70] & 62 [70] & 62 [70] \\
Fluid max.\ skewness & 1.31 [3.27] & 1.18 [2.64] & 1.27 [3.27] \\
Solid max.\ non-orthogonality (\si{\degree}) & 67 [75] & 68 [82] & 68 [82] \\
Solid max.\ skewness & 1.99 [3.42] & 1.38 [3.54] & 1.74 [3.54] \\
\addlinespace
\multicolumn{4}{l}{\textit{Surface, extrusion and lumen volume changes (\%)}} \\
Points moved: smoothing & 99.7 [99.0, 99.9] & 98.7 [97.4, 99.9] & 99.6 [97.4, 99.9] \\
Points moved: saddle rounding & 13.0 [9.0, 20.2] & 11.5 [4.8, 18.1] & 12.0 [4.8, 20.2] \\
Wall points: normal smoothed & 0.7 [0.0, 2.4] & 0.5 [0.0, 9.2] & 0.7 [0.0, 9.2] \\
Volume change: smoothing & 0.52 [0.22, 0.62] & 0.58 [0.42, 0.61] & 0.55 [0.22, 0.62] \\
Volume change: saddle rounding & 0.11 [$-$0.17, 0.39] & 0.11 [$-$0.32, 0.30] & 0.11 [$-$0.32, 0.39] \\
Volume change: total & 0.61 [0.13, 0.98] & 0.70 [0.12, 0.87] & 0.68 [0.12, 0.98] \\
\bottomrule
\end{tabular}
\end{table}


The automated meshing workflow was applied without manual intervention to all \MeshGeometriesAll{} geometries with $\varepsilon = 8$ (\cref{tab:meshing}).
It produced fluid and solid volume meshes for \MeshMeshedHuman{} of the \MeshGeometriesHuman{} human and \MeshMeshedPorcine{} of the \MeshGeometriesPorcine{} porcine geometries.
The solid mesh passed the standard \OpenFOAM{} mesh checks for \MeshCleanHuman{} human (\MeshCleanPctHuman\%) and \MeshCleanPorcine{} porcine (\MeshCleanPctPorcine\%) geometries, \MeshCleanAll{} in total (\MeshCleanPctAll\%).
Preparing and meshing a geometry took a median of \MeshRuntimeMedianMin~min on a single core, and at most \MeshRuntimeMaxMin~min.
\todo{State which geometries are carried forward to the simulations, once the mesh-acceptance rule of \cref{sec:Meshing} is decided.}

\paragraph{Failures before volume meshing}
No volume meshes were produced for \MeshFailedAll{} geometries (\MeshFailedHuman{} human, \MeshFailedPorcine{} porcine), for two reasons.
In \MeshOutletClosedAll{} (\MeshOutletClosedHuman{} human, \MeshOutletClosedPorcine{} porcine) a small outlet was closed in the fluid mesh, and the workflow stopped rather than lose its boundary condition.
Branches whose radius is below about \SI{1}{\milli\metre} along much of their length were affected as the branch is disconnected upstream of the local refinement at the outlet; the human trees, with more and finer branches (median of \MeshOpeningsMedianHuman{} openings against \MeshOpeningsMedianPorcine{}), were affected more often.
Further mesh refinement or larger refinement spheres would circumvent this but was not considered due to computational constraints.
The remaining \MeshScanDefectsAll{} failures (\MeshScanDefectsHuman{} human, \MeshScanDefectsPorcine{} porcine) were caused by defects in the segmented surfaces that no setting of the workflow can correct: two openings touching (\MeshTouchingOpeningsAll), two parts of the wall touching along an edge (\MeshNonManifoldWallAll), and a patch of wall duplicated and overlapping its neighbours (\MeshDuplicatedPatchAll).
These surfaces require repair of the segmentation.

\paragraph{Solid mesh quality}
Of the \MeshMeshedAll{} geometries that were meshed, \MeshOrientationAll{} solid meshes (\MeshOrientationHuman{} human, \MeshOrientationPorcine{} porcine) had incorrectly oriented faces or isolated highly skewed faces but no inverted cells, and \MeshNegativeAll{} (\MeshNegativeHuman{} human, \MeshNegativePorcine{} porcine) contained cells of negative volume, \MeshNegativeCellsTotal{} in total and at most \MeshNegativeCellsMax{} in a single mesh.
The meshes that passed the checks had a median of $\FluidCellsMedianHuman\times10^{3}$ fluid and $\SolidCellsMedianHuman\times10^{3}$ solid cells for the human geometries, and $\FluidCellsMedianPorcine\times10^{3}$ and $\SolidCellsMedianPorcine\times10^{3}$ for the porcine geometries, the difference reflecting the larger, more branched human trees.
Their maximum non-orthogonality did not exceed \FluidNonOrthoMaxAll\si{\degree} in the fluid or \SolidNonOrthoMaxAll\si{\degree} in the solid, and their maximum skewness did not exceed \FluidSkewMaxAll{} and \SolidSkewMaxAll{}, respectively.

\paragraph{Changes to the segmented surface}
The windowed-sinc smoothing moved almost every surface point (median \TaubinMovedMedianAll\%, a point counting as moved when displaced by more than \SI{0.02}{\milli\metre}), with a median of the mean displacement of \TaubinDispMedianAll~mm.
It increased the lumen volume by a median of \VolTaubinMedianAll\% (range \VolTaubinMinAll--\VolTaubinMaxAll\%).
The saddle rounding moved a median of \SaddleMovedMedianAll\% of the surface points (range \SaddleMovedMinAll--\SaddleMovedMaxAll\%), by a median of \SaddleDispMedianAll~mm where it acted and at most \SaddleDispMaxAll~mm, and changed the lumen volume by a median of \VolSaddleMedianAll\% (range \VolSaddleMinAll{} to \VolSaddleMaxAll\%).
Altogether the lumen volume changed by a median of \VolTotalMedianAll\% (range \VolTotalMinAll--\VolTotalMaxAll\%), with no systematic difference between the species (\cref{tab:meshing}).
The adaptive smoothing of the extrusion normals changed the normals of a median of \NormalSmoothedMedianAll\% of the wall points (at most \NormalSmoothedMaxAll\%); every other point was extruded along its own normal.

\subsection{Prestress}\label{sec:res-morph}
\todo{Discuss the results of the prestress and the differences between human and porcine}
\subsection{Flow and pressure}\label{sec:res-flow}
\todo{Waveforms, flow splits, pressure drop}
\subsection{Wall shear stress}\label{sec:res-wss}
\todo{TAWSS / OSI maps and distributions}
\subsection{Wall mechanics}\label{sec:res-wall}
\todo{Displacement, distensibility, wall stress}
